\documentclass[12pt]{article}
\usepackage[T1]{fontenc}
\usepackage[utf8]{inputenc}
\usepackage{lmodern}
\usepackage{textcomp}
\usepackage{enumitem}
\usepackage{geometry}
\geometry{margin=1in}
\begin{document}
\begin{center}
Answer Key - Health Sciences - Undergraduate Level Assessment 6 \
\small (Answers are ordered exactly as in the question paper.)
\end{center}

\section*{Multiple Choice}
\begin{enumerate}[label=\arabic*.]
\item \textbf{Answer:} D
\\ \textit{Options:} A) veins draining the scalp.; B) veins draining the eye.; C) the pterygoid venous plexus.; D) All of the above areas.
\\ \textit{Note:} Emissary veins serve as conduits that connect the intracranial venous sinuses to various regions, including the scalp, face, and diploë of the skull, thereby allowing for venous blood drainage from multiple areas.
\item \textbf{Answer:} A
\\ \textit{Options:} A) Take fewer medications; B) Be able to avoid certain diseases; C) Pay far less for pharmaceuticals; D) Pick and choose the disease(s) they want
\\ \textit{Note:} The correct answer is based on the understanding that identifying links between diseases can lead to more effective preventive measures and integrated treatment approaches, potentially reducing the need for multiple medications among older adults.
\item \textbf{Answer:} B
\\ \textit{Options:} A) erotica has no effect on sexual behavior; B) erotica may lead to increases in sexual arousal for men and women; C) nonaggressive erotica leads to massive increases in aggressive behavior; D) erotica decreases sexual arousal and sexual interest
\\ \textit{Note:} Research indicates that exposure to erotica, which is material that viewers find acceptable, can stimulate sexual arousal in both men and women, highlighting its potential psychological effects on sexual motivation and desire.
\item \textbf{Answer:} B
\\ \textit{Options:} A) The Baby Boom generation; B) The steadily increasing birth rate; C) Increases in life expectancy at birth; D) Increases in life expectancy at a specified age
\\ \textit{Note:} The steadily increasing birth rate is not a major reason for the increase in numbers of older adults; rather, the aging population is primarily attributed to longer life expectancies and declining birth rates, which lead to a higher proportion of older individuals in the population.
\item \textbf{Answer:} A
\\ \textit{Options:} A) That place holds many memories; B) It's much harder to move when you're older; C) Young adults pay less attention to place then do older adults; D) That place represents transcendence
\\ \textit{Note:} Older adults often have a deeper emotional connection to places due to the accumulation of life experiences and memories associated with those locations, which can foster a sense of attachment.
\end{enumerate}

\section*{True / False}
\begin{enumerate}[label=\arabic*.]
\setcounter{enumi}{5}
\item \textbf{Answer:} False
\\ \textit{Options:} A) True; B) False
\\ \textit{Note:} The labia majora are homologous to the scrotum in males, while the glans of the penis is homologous to the clitoris, making the original statement false.
\item \textbf{Answer:} False
\\ \textit{Options:} A) True; B) False
\\ \textit{Note:} Parasympathetic preganglionic nerves primarily exit the central nervous system through cranial nerves III, VII, IX, and X, not the fifth cranial nerve (trigeminal nerve), which is primarily involved in sensory functions and does not carry parasympathetic fibers.
\item \textbf{Answer:} True
\\ \textit{Options:} A) True; B) False
\\ \textit{Note:} The Philadelphia chromosome is formed when a piece of chromosome 22 is translocated to chromosome 9, resulting in the fusion of the BCR gene on chromosome 22 with the ABL gene on chromosome 9, which is characteristic of chronic myeloid leukemia.
\item \textbf{Answer:} True
\\ \textit{Options:} A) True; B) False
\\ \textit{Note:} The bones of the viscerocranium, which form the facial skeleton, develop through intramembranous ossification, allowing them to grow in size and shape in coordination with overall body growth, thus following the somatic growth pattern.
\item \textbf{Answer:} True
\\ \textit{Options:} A) True; B) False
\\ \textit{Note:} The statement is true because blood vessels can experience critical alterations, such as atherosclerosis or hypertension, which significantly increase the risk of cardiovascular events like heart attacks and strokes.
\end{enumerate}

\section*{Long Form Response}
\begin{enumerate}[label=\arabic*.]
\setcounter{enumi}{10}
\item \textbf{Answer:} Secondary sex characteristics play a crucial role in the sexual differentiation of males and females, influencing both biological and societal perceptions of gender. In males, traits such as facial hair, deeper voice, and increased muscle mass are linked to higher levels of testosterone, while females typically exhibit features like breast development and wider hips due to estrogen. These characteristics not only serve biological functions in reproduction but also shape societal views on masculinity and femininity. For instance, the presence of these traits can affect how individuals are perceived in terms of their roles, capabilities, and desirability in various contexts. Thus, secondary sex characteristics are significant as they contribute to both the physiological understanding of gender differences and the cultural constructs surrounding gender identity and expression.
\\ \textit{Note:} correct because it effectively outlines how secondary sex characteristics, influenced by hormones like testosterone and estrogen, contribute to both biological reproductive roles and societal gender perceptions by highlighting the distinct physical traits associated with males and females.
\item \textbf{Answer:} Hirschsprung disease exhibits a notable epidemiological trend in its gender distribution, as it is more common in girls than in boys, contrary to the traditional perception that it predominantly affects males. This condition, characterized by the absence of ganglion cells in the intestinal tract, leads to severe constipation and bowel obstruction. The misconception surrounding the prevalence may stem from the historical focus on male cases, which can skew public understanding and awareness of the disease. Additionally, misunderstanding the gender distribution can impact diagnosis and treatment approaches, leading to delays in care for affected individuals. Recognizing that Hirschsprung disease is more prevalent in females is crucial for improving health outcomes and ensuring timely intervention.
\\ \textit{Note:} correct because it addresses the unexpected gender distribution of Hirschsprung disease, noting its higher prevalence in girls, and highlights the implications of misconceptions based on historical biases that may lead to a misunderstanding of the condition's true epidemiology.
\end{enumerate}

\vfill
\noindent\textit{End of Answer Key}
\end{document}